\chapter{File types you will meet}
\label{app:filetypes}

ProPrep writes many files into your project folder. You never need to open most
of them, but it helps to recognize their names. This table lists the file types
that appeared in the test projects made for this manual, grouped by what they are
for. The ending of a file name (after the last dot) is called its
\term{extension}; it tells you what kind of file it is.

\begin{tip}
Most of these files are plain text. You can open them with any text editor
(TextEdit on a Mac, Notepad on Windows) to have a look. Do not change them by
hand unless a chapter tells you to.
\end{tip}

\section{Structures}

\begin{longtable}{@{}>{\ttfamily\bfseries}p{0.14\linewidth}>{\raggedright\arraybackslash}p{0.80\linewidth}@{}}
\toprule
\normalfont\bfseries Extension & \textbf{What it is} \\
\midrule
\endhead
.pdb & A protein structure in the classic Protein Data Bank format: one line per
  atom, with its name, residue, chain and position. Most ProPrep steps read and
  write this format. Each stage of preparation leaves its own copy, for example
  \file{filtered_structure.pdb} or \file{repaired_structure.pdb}
  (Chapter~\ref{ch:projectfolder}). \\
.cif & The newer mmCIF structure format. AlphaFold and AlphaFill models arrive
  as \file{.cif} files. Some ProPrep file browsers list only \file{.pdb} files
  (Chapter~\ref{ch:structure}). \\
.pdbqt & A structure prepared for docking with AutoDock Vina: a PDB file with
  partial charges and atom types added (Chapter~\ref{ch:docking}). \\
.sdf, .mol2 & Small-molecule files, used for ligands and docking poses.
  \file{.mol2} files also carry atom types and charges for parameterization
  (Chapter~\ref{ch:forcefields}). \\
.smi & A ligand written as a SMILES text string (Chapter~\ref{ch:docking}). \\
.xyz & A plain list of atom positions. ORCA writes one after a QM/MM run
  (Chapter~\ref{ch:qmmm}); docking writes a small one that describes the search box. \\
\bottomrule
\end{longtable}

\section{Simulation input}

\begin{longtable}{@{}>{\ttfamily\bfseries}p{0.14\linewidth}>{\raggedright\arraybackslash}p{0.80\linewidth}@{}}
\toprule
\normalfont\bfseries Extension & \textbf{What it is} \\
\midrule
\endhead
.prmtop, .parm7 & The \term{topology}: every atom, its charge and type, and
  every bond, angle and force constant. Built by the Topology Generator
  (Chapter~\ref{ch:topology}). A healthy one is usually several megabytes;
  a file of 0 bytes means the build failed. \\
.rst7 & The starting \term{coordinates} (and, later, velocities) that go with
  the topology. Always used as a pair with a \file{.prmtop} file. \\
.in (tLEaP) & The instructions ProPrep writes for \texttt{tleap}, the Amber
  program that builds the topology, for example \file{1UBQ_tleap.in}. You can
  view and edit it inside ProPrep before it runs. \\
.mdin & Settings for one simulation step (minimization, heating, production):
  how long, what temperature, which restraints (Chapter~\ref{ch:runningmd}). \\
.sh & Small scripts that run a simulation, locally or on a computing cluster
  (SLURM scripts). \\
.disang & Distance and angle restraints for a simulation (Chapter~\ref{ch:runningmd}). \\
.cpin & Settings for constant-pH simulations (Chapter~\ref{ch:protonation}). \\
.frcmod & Extra force-field parameters: bond, angle and twist strengths that the
  standard force field does not contain (Chapter~\ref{ch:forcefields}). \\
.lib & A residue library: the atoms, types and charges of a non-standard
  residue, such as a heme or a ligand. \\
.gjf, .com & Input files for Gaussian, the quantum-chemistry program used to
  derive new charges and for QM/MM (Chapters~\ref{ch:forcefields} and~\ref{ch:qmmm}). \\
.inp & Input files for ORCA, another quantum-chemistry program (Chapter~\ref{ch:qmmm}). \\
.gpf, .fld, .map & Grid files for AutoDock 4 scoring in docking (Chapter~\ref{ch:docking}). \\
\bottomrule
\end{longtable}

\section{Results and records}

\begin{longtable}{@{}>{\ttfamily\bfseries}p{0.14\linewidth}>{\raggedright\arraybackslash}p{0.80\linewidth}@{}}
\toprule
\normalfont\bfseries Extension & \textbf{What it is} \\
\midrule
\endhead
.nc & A \term{trajectory}: the protein's positions saved at intervals during a
  simulation, like the frames of a film (Chapter~\ref{ch:analysis}). \\
.mdout & The running log of a simulation: energies, temperature and pressure
  at each reporting interval. \\
.csv, .tsv & Tables you can open in Excel or any spreadsheet program: analysis
  results, docking scores, batch summaries. \\
.json & Structured text files ProPrep uses to remember things: the
  \term{session file} (\file{proprep_session_*.json}), redox-site definitions,
  workspace state, batch manifests (Chapters~\ref{ch:sessions} and~\ref{ch:redox}). \\
.log, .out & Logs from ProPrep and the programs it runs. When something fails,
  the real reason is often in a log, for example \file{leap.log} for topology
  building (Chapter~\ref{ch:troubleshooting}). \\
.pka & PROPKA's table of predicted pKa values (Chapter~\ref{ch:protonation}). \\
.ccp4 & Electron-density maps downloaded for structure validation
  (Chapter~\ref{ch:structure}). \\
.xml & BLAST search results exported from the Homology Searcher
  (Chapter~\ref{ch:compare}). \\
.txt & Plain-text reports and summaries. \\
\bottomrule
\end{longtable}
